\begin{lemma}
Let $H'$ and $H$ be finite-edge-indexed multihypergraphs on the same vertex
type, and suppose $H'$ is a sub-hypergraph of $H$ in the sense of
\noderef{SubhypergraphOf}. Then every uniformity parameter and every
simplicity parameter of $H$ is inherited by $H'$:
for every natural number $k$, if $H$ is $k$-uniform then $H'$ is
$k$-uniform; for every natural number $t$, if $H$ is $t$-simple then $H'$
is $t$-simple. Moreover, for every vertex $v$,
$\deg_{H'}(v)\le\deg_H(v)$.
\end{lemma}
\begin{proof}
By \noderef{SubhypergraphOf}, choose an injective edge-index map $i:F\to E$
with $H'.\mathrm{edge}(f)=H.\mathrm{edge}(i(f))$ for every $f$.
If every ambient edge has cardinality $k$, the same equality shows that
every child edge has cardinality $k$, proving the assertion about
\noderef{UniformHypergraph}.
For distinct child indices $f,g$, injectivity gives $i(f)\ne i(g)$.
The intersection of their vertex sets is exactly the intersection of the
ambient edges indexed by $i(f),i(g)$, which has cardinality at most $t$
by \noderef{TSimpleHypergraph}. This proves the simplicity assertion.
Finally, fix $v$. The map $i$ sends every child index incident with $v$
to an ambient index incident with $v$, by the edge equality, and remains
injective on this subset. The finite child incidence set therefore has
cardinality at most that of the ambient incidence set. These two
cardinalities are exactly the degrees by \noderef{HypergraphDegree}.
\end{proof}
